% Generated from locale/tr/content/first-order-logic/sequent-calculus/soundness-identity.tex; do not hand-edit.


\olfileid[tr]{fol}{seq}{sid}

\olsection{\usetoken{S}{identity} ile Sağlamlık}

\begin{prop}
Özdeşlik için başlangıç sekantları ve kurallar eklenmiş $\Log{LK}$ sağlamdır.
\end{prop}

\begin{proof}
${}\Sequent\eq[t][t]$ biçimindeki başlangıç sekantları geçerlidir; çünkü her
!!{structure}~$\Struct M$ için $\Sat{M}{\eq[t][t]}$. ($t$ teriminin kapalı,
yani değişken içermeyen bir terim olduğunu varsaydığımıza göre değişken
atamaları önemsizdir.)

Bir !!a{derivation} içindeki son çıkarımın $=$ kuralı olduğunu varsayalım. O zaman
öncül $\eq[t_1][t_2],\Gamma\Sequent\Delta,!A(t_1)$, sonuç ise
$\eq[t_1][t_2],\Gamma\Sequent\Delta,!A(t_2)$ olur. Bir
!!a{structure}~$\Struct M$ ele alalım. Sonucun geçerli olduğunu göstermemiz
gerekir; yani $\Sat{M}{\eq[t_1][t_2]}$ ve $\Sat{M}{\Gamma}$ ise ya bir
$!C\in\Delta$ için $\Sat{M}{!C}$ ya da $\Sat{M}{!A(t_2)}$.

Tümevarım hipotezine göre öncül geçerlidir. Bu, $\Sat{M}{\eq[t_1][t_2]}$
ve $\Sat{M}{\Gamma}$ ise ya (a) bir $!C\in\Delta$ için $\Sat{M}{!C}$ ya da
(b) $\Sat{M}{!A(t_1)}$ olduğu anlamına gelir. (a) durumunda işimiz bitmiştir.
(b) durumunu ele alalım. $s(x)=\Value{t_1}{M}$ olan bir $s$ değişken ataması
olsun. \olref[syn][ass]{prop:sentence-sat-true} uyarınca
$\Sat{M}{!A(t_1)}[s]$. $\varAssign{s}{s}{x}$ olduğundan
\olref[syn][ext]{prop:ext-formulas} uyarınca $\Sat{M}{!A(x)}[s]$.
$\Sat{M}{\eq[t_1][t_2]}$ olduğu için $\Value{t_1}{M}=\Value{t_2}{M}$ ve
dolayısıyla $s(x)=\Value{t_2}{M}$. \olref[syn][ext]{prop:ext-formulas}
önermesini yeniden uygulayınca $\Sat{M}{!A(t_2)}[s]$ elde ederiz.
\olref[syn][ass]{prop:sentence-sat-true} uyarınca $\Sat{M}{!A(t_2)}$.
\end{proof}

